\uuid{2tWw}
\exo7id{3628}
\titre{exo7 3628}
\auteur{quercia}
\organisation{exo7}
\datecreate{2010-03-10}
\isIndication{false}
\isCorrection{false}
\chapitre{Endomorphisme particulier}
\sousChapitre{Autre}
\cpge{mpsi-mp2i.espaces-vectoriels.a, pcsi.espaces-vectoriels.a, ptsi.geometrie-espace.c, mpsi-mp2i.espaces-vectoriels.c, pcsi.espaces-vectoriels.c, ptsi.espaces-vectoriels.c, mpsi-mp2i.dimension-finie.a, pcsi.dimension-finie.a, ptsi.dimension-finie.a, mpsi-mp2i.dimension-finie.b, pcsi.dimension-finie.b, ptsi.dimension-finie.b, bcpst1.espace-vectoriel-kn.a, bcpst1.espace-vectoriel-kn.b, bcpst2.espaces-vectoriels.a, bcpst2.espaces-vectoriels.b}
\module{Algèbre}
\niveau{L2}
\difficulte{}

\contenu{
\texte{
Soit $E =  K_n[X]$. Montrer que la famille ${\cal F} = (f_0,\dots,f_n)$ est
une base de $E^*$ et donner la base duale lorsque~\dots
}
\begin{enumerate}
    \item \question{$f_i(P) = P(x_i)$ où $x_0,\dots,x_n$ sont des scalaires distincts.}
    \item \question{$f_i(P) = P^{(i)}(0)$.}
    \item \question{$f_i(P) = P^{(i)}(x_i)$ où $x_0,\dots,x_n$ sont des scalaires quelconques.
    (Ne pas chercher la base duale pour cet exemple)}
\end{enumerate}
}
